\section{Método}
\label{sec:metodo}

Descreva o desenho da pesquisa de modo que outra pessoa possa compreender e,
quando aplicável, reproduzir a análise. Informe as fontes de dados, hipóteses,
variáveis, critérios de seleção, ferramentas, procedimentos e limitações.

Uma organização possível é:

\begin{enumerate}[leftmargin=*]
  \item definir o objeto e a unidade de análise;
  \item selecionar ou coletar os dados;
  \item aplicar o procedimento de análise;
  \item verificar a sensibilidade e interpretar os resultados.
\end{enumerate}

Se utilizar uma equação, defina cada símbolo logo após apresentá-la. Exemplo:

\begin{equation}
  C = \frac{N}{T},
  \label{eq:capacidade}
\end{equation}

em que $C$ representa a capacidade observada, $N$ o número de operações e $T$
o intervalo de tempo analisado. Substitua este exemplo pelo modelo do artigo.
